\begin{afterword}
\summaryhead

The second memoir post describes one belief of the young Yudkowsky as both a best and a worst mistake. Having seen AI researchers define intelligence loosely and then build cheap programs that met the definition, the young author refused to define it. Having seen neat mathematics confine AI to toy systems, the young author disbelieved in any mathematics of intelligence, except Bayes's theorem. The author now calls both overreactions to other people's errors, and traces them to reading too few, too popular AI books and to writing the field off.

The belief was a best mistake because of what it led to: wide study of the cognitive sciences, no stopping at the first bright idea, and long study before proposing any system. The moral is that what you end up doing ``screens off'' the clever reason for doing it. It was a worst mistake because informal reasoning was held to low standards. Such reasoning is needed on the way to knowledge, but must not be trusted, and the young author trusted it.

\responsehead

The post is memoir, and it is judged as memoir. The author's account of the young author's reasons and feelings is taken as given. What can be checked is the account of the young author's views against the writings of the time, and most of it is accurate.

``The Plan to Singularity'' (1999--2000) confirms the low view of AI. It says the needed talent ``is more likely to reside in the field now called `cognitive science' than in the field called AI,'' and that ``AI remains crippled.'' It also shows the neuroanatomical picture of minds the post describes: a module for source code is called a ``codic cortex,'' ``by analogy to the human `visual cortex'.''

One description is stronger than the record. The post says the young author ``would have been loath to put any definition on `intelligence' at all,'' and two paragraphs later, ``I refused to define it at all.'' The archived texts do define intelligence, informally. ``Staring into the Singularity'' has a section headed ``The Definition of Smartness,'' and the Plan says ``Intelligence is whatever lets you model, predict, and manipulate reality.'' Neither is mathematical, so the report that the young author refused a mathematical definition is accurate, and the post's later wording, that the young author thought it ``okay \ldots\ not to precisely define'' intelligence, fits the record.

The post is candid about luck. It calls many of the author's successes right actions for wrong reasons, and turns the selection effect against its own record: readers hear only from an author who did not reach a dead end.

The post leaves the math of intelligence it came to accept to a later post, ``My Bayesian Enlightenment.'' Formal work of that kind existed by then: Legg and Hutter gave a formal definition of machine intelligence in 2007.

The post ends with ``To be continued'' before saying which informal conclusion the young author trusted. ``A Prodigy of Refutation,'' two posts on, returns to the belief that a superintelligence would know what is right.

In short: a memoir, candid about luck, whose account of the young author's views on AI is supported by the archived writings, though those writings did define intelligence informally.
\end{afterword}
